\subsubsection{Generación de casos de prueba}

Los estudiantes deberán implementar un generador de casos de prueba en Python llamado \texttt{testcases\_generator.py}, ubicado en:

\begin{center}
    \texttt{code/implementation/scripts/testcases\_generator.py}
\end{center}

\textbf{Requisitos del generador:}
\begin{itemize}
    \item Debe generar archivos de entrada válidos para el problema \textbf{AniMarathon}.
    \item Cada archivo de entrada debe llamarse:
    \[
        \texttt{testcases\_\{n\}\_\{i\}.txt}
    \]
    donde:
    \begin{itemize}
        \item \texttt{\{n\}} representa la cantidad de animes del caso de prueba.
        \item \texttt{\{i\}} es un identificador único para diferenciar múltiples casos con la misma cantidad de animes.
    \end{itemize}

    \item Formato del archivo de entrada:
    \begin{itemize}
        \item La primera línea contiene tres enteros $n$, $M$ y $E$, donde:
        \begin{itemize}
            \item $n$ es la cantidad de animes disponibles.
            \item $M$ es la cantidad máxima de minutos disponibles.
            \item $E$ es la energía máxima disponible.
        \end{itemize}

        \item Para cada anime $i$, se entrega una línea con:
        \[
        \texttt{anime}_i \ q_i \ b_i
        \]
        donde:
        \begin{itemize}
            \item $\texttt{anime}_i$ es el nombre del anime, representado por un string sin espacios.
            \item $q_i$ es la cantidad de capítulos del anime.
            \item $b_i$ es el bono de satisfacción por completar todos los capítulos del anime.
        \end{itemize}

        \item Luego, para cada uno de los $q_i$ capítulos del anime $i$, se entrega una línea con:
        \[
        t_{i,j} \ c_{i,j} \ v_{i,j}
        \]
        donde:
        \begin{itemize}
            \item $t_{i,j}$ es la duración del capítulo en minutos.
            \item $c_{i,j}$ es el costo de energía del capítulo.
            \item $v_{i,j}$ es la satisfacción obtenida al ver el capítulo.
        \end{itemize}
    \end{itemize}

    \item El generador debe permitir crear múltiples casos de prueba de manera automática, variando la cantidad de animes, cantidad de capítulos, minutos disponibles, energía disponible, puntajes de satisfacción y bonos de completación.
\end{itemize}

\begin{mdframed}
    Los archivos generados serán utilizados por los algoritmos implementados en C++ para probar la corrección y eficiencia de las soluciones de fuerza bruta, greedy y programación dinámica.
\end{mdframed}

\subsubsection{Tipos de casos de prueba esperados}

El generador debe ser capaz de producir, al menos, los siguientes tipos de casos:

\begin{enumerate}
    \item \textbf{Casos pequeños:} utilizados para validar la correctitud del algoritmo de fuerza bruta y compararlo con programación dinámica. Se recomiendan casos con pocos animes y pocos capítulos, por ejemplo:
    \[
        n \in \{3, 5, 8\}
    \]
    con una cantidad total de capítulos pequeña.

    \item \textbf{Casos medianos:} utilizados para comparar las heurísticas greedy con la solución óptima obtenida por programación dinámica. Se recomiendan valores como:
    \[
        n \in \{20, 40, 80\}
    \]

    \item \textbf{Casos grandes:} utilizados para medir rendimiento de los algoritmos eficientes. Se recomiendan valores como:
    \[
        n \in \{100, 150, 200\}
    \]
    En estos casos no se espera que fuerza bruta sea ejecutable en tiempos razonables.
\end{enumerate}
    Se recomienda usar el tipo de dato long long para prevenir errores de redondeo

\subsubsection{Ejemplo de caso de prueba válido}

El siguiente archivo corresponde a un caso de prueba válido para el problema:

\begin{verbatim}
4 240 90
shonen_quest 3 25
45 18 40
50 22 45
55 25 60
romcom_days 2 15
30 10 25
35 12 30
mecha_nova 2 50
60 30 65
70 35 75
slice_cafe 1 5
25 8 18
\end{verbatim}

\subsubsection{Validaciones esperadas}

Los casos de prueba generados deben permitir verificar que:

\begin{itemize}
    \item Los algoritmos no seleccionen capítulos aislados sin seleccionar los capítulos anteriores del mismo anime.
    \item Los algoritmos respeten el límite de minutos disponibles $M$.
    \item Los algoritmos respeten el límite de energía disponible $E$.
    \item Los bonos de completación se sumen únicamente cuando se seleccionan todos los capítulos de un anime.
    \item El algoritmo de fuerza bruta y el de programación dinámica entreguen el mismo resultado en casos pequeños.
    \item Las heurísticas greedy entreguen soluciones válidas, aunque no necesariamente óptimas.
    \item La programación dinámica pueda resolver casos medianos y grandes dentro de tiempos razonables.
    \item Los gráficos permitan observar diferencias claras entre los enfoques implementados.
\end{itemize}